% Section 5: Triadic characteristic inference
\label{sec:triadic-inference-body}

\S\ref{sec:observers} fixed what an observer is and what a characteristic costs.
This section reads one finite characteristic algebra in three ways at once and asks
what the three readings share.  The answer has two parts: a variance signature that is
stable under every calibration performed so far, and an inference procedure that is
nothing more than the intersection of three fibres.  Both are developed on the
affine monoid of Lemma~\ref{lem:affine-product} and then calibrated on the square and
the cube in \S\ref{sec:multiplicity}.

\subsection{The three readings}

\begin{definition}[Typed readings]\label{def:typed-readings}
Let $\cA$ be a class of processes closed under a composition $\odot$ read in execution
order.  A \emph{typed reading family} on $\cA$ consists of
\begin{enumerate}
\item a \emph{temporal} reading $\rho_t$, covariant in execution, taking values in
  endomorphisms of the state space;
\item a \emph{spatial} reading $\rho_X$, contravariant, taking values in inverse-image
  maps on a space of admissible regions;
\item a \emph{constructive} reading $\rho_\cK$, covariant modulo a declared
  canonicalisation $\equiv_{\mathrm{can}}$, taking values in construction data.
\end{enumerate}
\end{definition}

The three readings have different variance by design: a process acts forward on states,
backward on regions, and forward again on constructions.  This asymmetry is not a
defect; it is what the phrase "three readings" is meant to encode.

\begin{proposition}[Variance signature]\label{prop:variance-signature}
Let $f_1$ be executed before $f_2$ and write $f_2\odot f_1$ for the composite.  Then
\[
  \rho_t(f_2\odot f_1)=\rho_t(f_2)\circ\rho_t(f_1),\qquad
  \rho_X(f_2\odot f_1)=\rho_X(f_1)\circ\rho_X(f_2),
\]
\[
  \rho_\cK(f_2\odot f_1)\equiv_{\mathrm{can}}\rho_\cK(f_2)\circ\rho_\cK(f_1).
\]
Schematically, writing $\mathfrak T$ for the triple,
\[
  \mathfrak T(f_2\odot f_1)=\bigl(T_2\circ T_1,\ X_1\circ X_2,\ K_2\circ K_1
  \bmod \mathrm{can}\bigr),
\]
so the signature is $({+},{-},{+})$ in the order $(t,X,\cK)$.
\end{proposition}

\begin{proof}
The temporal and constructive statements are functoriality of $\rho_t$ and of
$\rho_\cK$ modulo canonicalisation.  The spatial statement is contravariance of the
inverse-image construction: a region transformed by the composite $f_1$ then $f_2$ is
$\rho_X(f_2)^{-1}\bigl(\rho_X(f_1)^{-1}(R)\bigr)=\bigl(\rho_X(f_1)\circ\rho_X(f_2)
\bigr)^{-1}(R)$.
\end{proof}

\begin{calibration}[Scope of the signature]\label{cal:variance-scope}
The signature of Proposition~\ref{prop:variance-signature} is established by hand
inside the finite affine model, and re-verified in the calibrations of
\S\ref{sec:multiplicity} and \S\ref{sec:classical}.  It is \emph{not} claimed as a
theorem about every conceivable three-domain model.  What changes from model to model
is the structure each reading must carry; the variance law has so far not changed.
\end{calibration}

\subsection{Inference as intersection of fibres}

\begin{proposition}[Triadic inference]\label{prop:triadic-inference}
For observations $o_t,o_X,o_\cK$ presented to the three readings, the inferred set of
candidates is
\[
  \cI(o_t,o_X,o_\cK)
  =\Fib_{t,Q}(o_t)\cap\Fib_{X,Q}(o_X)\cap\Fib_{\cK,Q}(o_\cK).
\]
Each reading contributes constraints and the constraints are intersected.
\end{proposition}

\begin{proof}
By Definition~\ref{def:observation-fibre}, a process lies in $\Fib_{D,Q}(o_D)$ exactly
when its $D$-reading is $Q$-indistinguishable from $o_D$.  A process satisfying all
three observations therefore lies in the intersection, and conversely.
\end{proof}

\begin{remark}[What "one algorithm" can and cannot mean]\label{rem:one-algorithm}
Proposition~\ref{prop:triadic-inference} says that the three readings share one
\emph{operation} in characteristic space.  It does not say that they share an untyped
decoder; the fibres live in different codomains and the intersection is taken after
transport.  Reading the proposition as a claim about a universal interpreter is a
mistake, and it is the mistake that Part~II is designed to prevent.
\end{remark}

The outcome of the intersection is reported through
Definition~\ref{def:outcome-discipline}.  In particular
Warning~\ref{warn:unknown-not-nonexistence} applies with force here: an empty
intersection is \emph{inconsistent} only if all three searches were exhausted within
$B_Q$.

\subsection{The deformed crossing}

Inference is only half the story; the other half is how readings move past one
another.  The affine model gives an exact law, and it is not a swap.

\begin{proposition}[Order reduction is not a swap]\label{prop:deformed-crossing}
Write $S$ for successor and $D$ for doubling.  Then
\[
  S;D\equiv D;S;S,\qquad\text{while}\qquad S;D\not\equiv D;S,
\]
where $\equiv$ denotes extensional equality of the induced value.  Moving $D$ leftward
across $S$ \emph{duplicates} the translation generator, and every pair has a canonical
representative $D^{k};S^{b}$.
\end{proposition}

\begin{proof}
By Lemma~\ref{lem:affine-product}, $S;D$ has parameters $(1,2)$ while $D;S$ has
$(1,1)$; and $D;S;S$ has $(1,2)$ as well.  Canonicality of $D^{k};S^{b}$ is the normal
form of the monoid.
\end{proof}

\begin{warning}[Extensional, not confluent]\label{warn:not-confluent}
Proposition~\ref{prop:deformed-crossing} is one exact extensional crossing law.  It is
\emph{not} a terminating and confluent rewriting theorem for the abstract
presentation, and this paper does not claim one.  The distinction matters because
Part~II proves genuine normalization theorems under linearity hypotheses, and
Proposition~\ref{prop:deformed-crossing} must not be mistaken for one of them.
\end{warning}

\subsection{The square map: loss, branching, and reverse transport}

The clearest model of what a reading can and cannot retain is the real square.

\begin{proposition}[Temporal loss is an orbit quotient]\label{prop:square-orbit}
Let $q(x)=x^{2}$ and $\sigma(x)=-x$.  Then $q\circ\sigma=q$, so $q$ factors as
\[
  \R\xrightarrow{\ \mathrm{quot}\ } \R/C_2\cong[0,\infty)
  \xrightarrow{\ \hookrightarrow\ } \R .
\]
The generic fibre is a single two-element $C_2$-orbit; $\{0\}$ is a fixed point, so the
quotient is non-free exactly at the origin.
\end{proposition}

\begin{proposition}[Spatial branching is the same quotient, read backwards]
\label{prop:square-covering}
On $\R^{\times}\to\R_{>0}$, $q$ is a two-sheeted covering and $\sigma$ exchanges the
sheets.  For $0<a<b$,
\[
  q^{-1}(a,b)=(-\sqrt b,-\sqrt a)\cup(\sqrt a,\sqrt b),
\]
while $q^{-1}(-1,1)=(-1,1)$ exhibits the sheets merging through the critical source
point.
\end{proposition}

\begin{proposition}[Reverse transport is relation-valued]\label{prop:relation-valued}
The backward reading $q^{\dagger}$ is $\varnothing$, $\{0\}$, or
$\{-\sqrt y,\sqrt y\}$ according as $y<0$, $y=0$, or $y>0$.  Hence a plain global
inverse function does not exist, although forward execution remains functional.
\end{proposition}

\begin{remark}
Propositions~\ref{prop:square-orbit}--\ref{prop:relation-valued} are the smallest
complete instance of the warning that contravariance forces relation-valued transports.
The spatial reading of a critical source is not a function, and any formalism that
demands one will either lose the singularity or invent a branch convention that is not
intrinsic to the data.
\end{remark}

\subsection{Scaling crosses; translation does not}

\begin{proposition}[Scaling crosses by deformation]\label{prop:scaling-crosses}
For $D_a(x)=ax$ and $Q(x)=x^{2}$,
\[
  D_a;Q\equiv Q;D_{a^{2}},
\]
with equal extensional value and different raw histories.
\end{proposition}

\begin{proof}
$q(D_a x)=a^{2}x^{2}$ and $D_{a^{2}}(q(x))=a^{2}x^{2}$.
\end{proof}

\begin{theorem}[No-go: translation has no affine crossing]\label{thm:translation-obstruction}
Let $S_b(x)=x+b$ with $b\neq0$.  There is no affine $A(y)=\alpha y+\beta$ with
\[
  S_b;Q\equiv Q;A .
\]
\end{theorem}

\begin{proof}
Every post-composite of $Q$ with an affine map has the form
$A(q(x))=\alpha x^{2}+\beta$, whose linear coefficient vanishes identically.  But
$q(S_b x)=x^{2}+2bx+b^{2}$ has linear coefficient $2b\neq0$ since the field has
characteristic $0$ and $b\neq0$.
\end{proof}

\begin{remark}[Scope of Theorem~\ref{thm:translation-obstruction}]
\label{rem:obstruction-scope}
The theorem forbids an \emph{affine} crossing.  It does not forbid a non-affine
expansion, and the next result exhibits exactly that.  Reading the obstruction as a
claim about all rewrites would overstate it; this is the sense in which the statement
is a boundary statement about the affine post-normal-form language, not about
reduction in general.
\end{remark}

\begin{proposition}[The word-to-graph boundary]\label{prop:word-to-graph}
Crossing a translation through a square forces a parallel decomposition with explicit
copy:
\[
  x\xrightarrow{\ \Delta_3\ } (x_1,x_2,x_3)
  \xrightarrow{\ Q\otimes D_{2b}\otimes C_{b^{2}}\ } (x_1^{2},\,2bx_2,\,b^{2})
  \xrightarrow{\ \mathrm{add}_3\ } x^{2}+2bx+b^{2}.
\]
Consequently general reduction in this setting is not a global ordering of two word
alphabets, but typed graph rewriting with declared copy and merge.
\end{proposition}

\begin{proof}
The right-hand side of the obstruction is a sum of three terms whose coefficients come
from three distinct uses of one source.  A word over the one-input alphabets $\{S_b\}$
and $\{Q\}$ cannot express three simultaneous uses of its single input; a graph with an
explicit fan-out can.
\end{proof}

\begin{calibration}[Constructive fibre is not a singleton]\label{cal:constructive-fibre}
$\chi_\cK(Q)=\chi_\cK(N;Q)=q$ while the raw histories and the construction data of $Q$
and $N;Q$ remain distinct.  Hence the constructive fibre is not a singleton even when
the temporal and spatial fibres are, and by
Proposition~\ref{prop:residual-principle} a usable characteristic must retain the
residual.
\end{calibration}

\subsection{Summary}

Three results in this section are load-bearing for what follows.

\begin{enumerate}
\item The variance signature $(+,-,+)$ is stable across the models tested
  (Proposition~\ref{prop:variance-signature}, Calibration~\ref{cal:variance-scope}).
\item Crossings are deformations, and translation does not cross an affine normal form
  at all (Propositions~\ref{prop:deformed-crossing}, \ref{prop:scaling-crosses},
  Theorem~\ref{thm:translation-obstruction}).
\item The consequence is a change of carrier: from words to typed graphs with explicit
  copy (Proposition~\ref{prop:word-to-graph}).  Part~II takes up precisely this
  carrier, where the copy discipline becomes a linearity discipline.
\end{enumerate}
